%----------------------------------------------------------------------------
\chapter{Design}
\label{sec:Design}
%----------------------------------------------------------------------------
% Visible marker for a decision whose original reason is not documented.
% Every \reasontodo must be resolved (or removed) before submission.
\providecommand{\reasontodo}[1]{\textcolor{red}{\textbf{[TODO: #1]}}}

This chapter describes the design of the automated rollback platform: what it has to achieve, how it is put together, and why each part looks the way it does. The chapter describes the final design; how it was built, and the problems met on the way, are the subject of \autoref{sec:Implementation}.

%----------------------------------------------------------------------------
\section{Requirements}
\label{sec:DesignRequirements}
%----------------------------------------------------------------------------
The goal of the platform is that a faulty release is detected from the behaviour of the new version under live traffic and reverted without a human. The requirements follow from this goal and from the constraints of the available infrastructure: two small virtual machines with 8\,GB of memory and 4 CPU cores each, reachable from outside only over SSH, and without a public domain name. \autoref{tab:DesignRequirements} lists them, together with the quantity that shows whether a requirement is met and the place where it is checked.

\begin{table}[ht]
	\footnotesize
	\centering
	\begin{tabular}{ l p{5.2cm} p{5.0cm} l }
		\toprule
		ID & Requirement & Target & Checked in \\
		\midrule
		R1 & A faulty release is aborted and all users are returned to the old version automatically. & no manual step; recovery within seconds of the detection & \autoref{sec:EvalTTD} \\ \addlinespace
		R2 & A fault visible in the request metrics reaches as few users as possible. & at most the first canary step (10\,\% of the users); 0\,\% where the fault can be caught on copied traffic & \autoref{sec:EvalImpact} \\ \addlinespace
		R3 & The fault is detected quickly. & within about 1--2 minutes of the first request to the new version, limited by a 10\,s scrape interval & \autoref{sec:EvalTTD} \\ \addlinespace
		R4 & A healthy release is not rolled back. & no aborted healthy release & \autoref{sec:EvalDetection} \\ \addlinespace
		R5 & The kubernetes platform fits on one small VM. & 8\,GB of memory per VM & \autoref{sec:ImplEnvironment} \\ \addlinespace
		R6 & An environment can be rebuilt from code. & one Ansible command takes a fresh VM to a synchronised cluster; the trusted CA survives a rebuild & \autoref{sec:ImplBootstrap} \\ \addlinespace
		R7 & The platform is private and its secrets are protected. & TLS from an own CA; no plaintext secret in Git & \autoref{sec:ImplNetworking} \\ \addlinespace
		R8 & Every release, and every experiment, is traceable and repeatable. & every release is a Git commit; a trial needs no manual step & \autoref{sec:EvalRunner} \\
		\bottomrule
	\end{tabular}
	\caption{Requirements of the platform.}
	\label{tab:DesignRequirements}
\end{table}

R1--R4 describe the rollback itself and pull in different directions. Exposing fewer users (R2) means less traffic on the new version, which makes its metrics noisier and either slower to judge (R3) or more prone to false alarms (R4). Finding a good balance between them, and measuring it, is the subject of the experiments in \autoref{sec:Evaluation}. R5--R8 constrain how the platform is built.

Two limits of the scope follow from the same goal. First, the platform protects against faulty \emph{releases}, not against the loss of a machine: the cluster is a single node, without high availability, and a lost VM is rebuilt from code (R6). Second, only faults that are part of the release are considered. The faults of the environment were left out by decision, such as a node under CPU pressure, to keep the experiments focused on the question of when a new version should be rolled back.

%----------------------------------------------------------------------------
\section{Architecture overview}
\label{sec:DesignOverview}
%----------------------------------------------------------------------------
\autoref{fig:DesignOverview} shows the system. It spans three places. The \emph{control machine}, runs the provisioning tools and the trial runner. A second VM, \texttt{prgdlv-auto}, hosts the Git server Gitea together with its container registry and its CI runner, which I set up as host services outside any cluster. The cluster itself runs on \texttt{prgdlv-test}. All three are connected by the NetBird overlay VPN, through which the platform and its user interfaces are reached by name.

\begin{figure}[ht]
	\centering
	\resizebox{\textwidth}{!}{%
	\begin{tikzpicture}[
		box/.style={draw, rounded corners=2pt, align=center, font=\small, inner sep=4pt, minimum height=1.0cm},
		comp/.style={box, fill=black!5, text width=3.2cm},
		lbl/.style={font=\scriptsize, align=center},
		arr/.style={-{Stealth[length=2mm]}, semithick}]
		% top row
		\node[box, text width=4.0cm] (ctrl) at (2.4,7.2) {\textbf{Control machine}\\Ansible, kubectl, Helm\\SOPS with the age key\\trial runner, browser};
		\node[box, text width=5.6cm] (auto) at (12.2,7.2) {\textbf{prgdlv-auto} (host services)\\Gitea: \texttt{rollback-platform}, \texttt{demo}\\container registry, CI runner};
		% cluster
		\node[comp] (gw)    at (2.2,3.9)  {ServiceLB and\\Istio ingress gateway};
		\node[comp] (demo)  at (2.2,1.3)  {namespace \texttt{demo}:\\stable, canary pods,\\load generator};
		\node[comp] (argo)  at (6.8,3.9)  {Argo CD};
		\node[comp] (ro)    at (6.8,1.3)  {Argo Rollouts};
		\node[comp] (cm)    at (11.6,3.9) {cert-manager, istiod};
		\node[comp] (prom)  at (11.6,1.3) {Prometheus};
		\node[comp] (dash)  at (16.4,1.3) {Grafana, Kiali};
		\coordinate (spacer) at (2.2,-0.05);
		\begin{scope}[on background layer]
			\node[draw, thick, rounded corners=4pt, fit=(gw)(demo)(argo)(ro)(cm)(prom)(dash)(spacer), inner sep=10pt, label={[font=\small\bfseries, anchor=south east]north east:prgdlv-test: single-node k3s cluster (dev)}] (cluster) {};
			\node[draw, dashed, rounded corners=6pt, fit=(ctrl)(auto)(cluster), inner sep=8pt, label={[font=\small, anchor=south west]north west:NetBird VPN}] {};
		\end{scope}
		% arrows
		\draw[arr] (ctrl) -- node[lbl, above]{\texttt{git push}} (auto);
		\draw[arr] (ctrl.south) -- node[lbl, left, pos=0.55]{HTTPS\\\texttt{*.dev.internal}} (gw.north);
		\draw[arr] (ctrl.south east) -- node[lbl, right, pos=0.35]{bootstrap\\(once)} (argo.north west);
		\draw[arr] (auto.south west) -- node[lbl, right, pos=0.45]{pull Git\\(polling)} (argo.north east);
		\draw[arr] (argo) -- node[lbl, right]{applies\\manifests} (ro);
		\draw[arr] (ro) -- node[lbl, above, sloped]{route weights} (gw);
		\draw[arr] (gw) -- node[lbl, left]{routes users} (demo);
		\draw[arr] (ro) -- node[lbl, above]{queries} (prom);
		\draw[arr] (dash) -- node[lbl, above]{queries} (prom);
		\draw[arr] (prom.south) -- ++(0,-0.7) -| node[lbl, above, pos=0.25]{scrapes sidecars and cAdvisor (pull)} (demo.south);
	\end{tikzpicture}}
	\caption{Overview of the platform. Images are pulled from the registry on \texttt{prgdlv-auto}; this is not drawn.}
	\label{fig:DesignOverview}
\end{figure}

%----------------------------------------------------------------------------
\subsection{Two layers: bootstrap and GitOps}
\label{sec:DesignLayers}
%----------------------------------------------------------------------------
The platform is set up in two layers with a strict boundary between them. The \emph{bootstrap layer}, a pair of Ansible playbooks run from the control machine, installs Kubernetes on the VM, creates the few objects that cannot come from Git, installs Argo CD, and hands over to it by creating a single root Application. From then on, the \emph{GitOps layer} owns the cluster: everything inside it is described in the Git repository and applied by Argo CD.

The boundary is placed there for two reasons. First, GitOps keeps the desired state declarative and versioned in Git, and an agent in the cluster pulls it and keeps reconciling the cluster with it~\cite{OpenGitOps}. This gives R8 for free: every change, including every release, is a commit that can be traced, reverted and cited. It also makes manual drift impossible to keep, because Argo CD's self-heal undoes it. Second, the objects that the bootstrap creates are exactly those that Git cannot hold. Argo CD cannot decrypt the encrypted secrets of the repository, so the secrets (the root CA, the Grafana administrator password and the Kiali signing key) are decrypted by Ansible and created as Kubernetes Secrets before Argo CD starts. The applications in Git refer to these Secrets by name only. After the root Application has been applied, Ansible is never needed again until the VM is rebuilt.

%----------------------------------------------------------------------------
\section{Choice of components}
\label{sec:DesignComponents}
%----------------------------------------------------------------------------
\autoref{tab:DesignComponents} summarises the components, the alternatives that were considered and the main reason for each choice. The paragraphs below give the reasons in more detail.

\begin{table}[ht]
	\footnotesize
	\centering
	\begin{tabular}{ p{2.2cm} p{2.5cm} p{3.2cm} p{4.9cm} }
		\toprule
		Concern & Choice & Alternatives & Main reason \\
		\midrule
		Kubernetes & k3s, single node & full Kubernetes distribution & small memory footprint on an 8\,GB VM \\ \addlinespace
		Load balancer & k3s ServiceLB & MetalLB & one VM per cluster; nothing to fail over to \\ \addlinespace
		Ingress, mesh & Istio (sidecar mode) & Traefik (bundled with k3s) & one component for ingress, weighted routing, mirroring and request metrics \\ \addlinespace
		TLS & cert-manager with an own CA & public CA with ACME & no public domain; private \texttt{.internal} names \\ \addlinespace
		Secrets & SOPS with age & Ansible Vault, Sealed Secrets, external secret store & readable diffs, per-person keys, usable outside Ansible, no extra server \\ \addlinespace
		Git, registry, CI & standalone Gitea on prgdlv-auto & Git server inside the cluster & no dependency loop: the cluster's configuration comes from Git \\ \addlinespace
		Access & NetBird VPN, \texttt{.internal} names & SSH to each VM & every service reachable by name from one network \\ \addlinespace
		Delivery controller & Argo Rollouts & Flagger & see the paragraph below \\ \addlinespace
		Metrics & standalone Prometheus & logs, traces & see \autoref{sec:MetricAnalysis} \\ \addlinespace
		Dashboards & Grafana and Kiali & Prometheus UI only & time series of a release; live mesh graph \\
		\bottomrule
	\end{tabular}
	\caption{Components of the platform and the reasons for choosing them.}
	\label{tab:DesignComponents}
\end{table}

\paragraph{Kubernetes distribution.} The clusters run k3s, a Kubernetes distribution packaged as a single binary that needs about half the memory of a standard installation~\cite{K3s}. Its footprint matters on an 8\,GB VM that also has to hold a service mesh, a metric store and two delivery controllers (R5). Besides the memory, three further reasons favoured k3s:
\begin{itemize}
	\item \emph{Full Kubernetes.} k3s is a fully compliant Kubernetes distribution~\cite{K3s}, so Argo CD, Argo Rollouts and Istio run on it as on any other cluster, and the results of the evaluation do not depend on a simplified cluster.
	\item \emph{Batteries included.} k3s brings the components that a single node needs: the containerd runtime, the Flannel network, CoreDNS, a provisioner for local volumes and the ServiceLB load balancer~\cite{K3s}. With a standard installation, each of these would have to be chosen and installed separately. Its bundled ingress controller, Traefik, is disabled, since Istio takes its place.
	\item \emph{Easy to rebuild.} The official Ansible collection of k3s installs a cluster from one playbook~\cite{K3sAnsible}, so an environment can be wiped and rebuilt from the repository (\autoref{sec:ImplBootstrap}).
\end{itemize}
Lighter local tools, such as kind or minikube, are made for short-lived clusters on a developer's machine rather than for a cluster that runs on a VM and is reached over a VPN. The price of k3s in this form is that a single server stores the cluster state in SQLite rather than etcd and has no high availability~\cite{K3s}; for the rollback logic, which runs above this layer, this makes no difference.

\paragraph{Load balancer and ingress.} k3s comes with ServiceLB, which implements Services of type \texttt{LoadBalancer} by opening host ports on the node~\cite{K3sServiceLB}. In effect, the node itself is the load balancer. The main alternative, MetalLB, gives each Service its own virtual IP address on the network, which would allow several Services to share port 443 and would move the address if a node failed. With one VM per cluster, there is no second node to move to, and only the Istio gateway needs ports 80 and 443, because every application is reached through it. MetalLB would therefore add complexity without benefit, and I kept ServiceLB. The ingress controller bundled with k3s, Traefik, is disabled: the Istio ingress gateway is the only entry point, since Istio is needed anyway (see below), and two ingress controllers would compete for the same ports.

\paragraph{Service mesh.} Istio is the core of the detection: its sidecar proxies split traffic between the stable and the canary version by weight, copy requests to the canary for mirroring, and record request metrics without any change to the application (\autoref{sec:MetricAnalysis}). Argo Rollouts supports Istio directly by rewriting the weights of a VirtualService~\cite{ArgoRolloutsIstio}. The mesh runs in sidecar mode, with sidecar injection enabled only in application namespaces and strict mutual TLS (\texttt{PeerAuthentication} in \texttt{STRICT} mode) in them, so that only proxies holding a mesh certificate can call the applications~\cite{IstioPeerAuth}. The design depends on the sidecar mode in two places: the checks read the request metrics that the new pod's own sidecar records (\autoref{sec:DesignChecks}), and the weighted routing and the mirroring of the canary are HTTP features of the proxies.

\paragraph{Why sidecar mode rather than ambient mode.} Istio offers a second data plane mode, the ambient mode, in which the application pods contain no proxy. A proxy on each node, the ztunnel, handles only layer~4: mutual TLS, authorization by workload identity and TCP metrics. The HTTP features, among them per-request routing, mirroring and the request metrics with status codes and latency, require an additional Envoy deployment, a waypoint proxy, in front of the workloads~\cite{IstioDataplaneModes,IstioWaypoint}. I nevertheless chose the sidecar mode, for four reasons:
\begin{itemize}
	\item \emph{Per-pod metrics.} Every pod's own proxy reports the success rate and the latency of the requests it served, which is exactly what is needed to judge one new version. In ambient mode, these metrics would be recorded by the waypoint, and every query would have to select the new version by workload labels instead of by the pod that Prometheus scraped.
	\item \emph{Maturity.} The sidecar mode is the older mode, and Istio supports it in more settings than the ambient mode, for example across multiple clusters and for virtual machines~\cite{IstioDataplaneModes}. A rollback platform should not depend on the newer features of its mesh.
\end{itemize}

\paragraph{TLS.} Without a public domain, a public certificate authority cannot issue certificates for the platform's host names, which live under the private suffix \texttt{.internal}. The environment therefore has its own root CA, from which cert-manager's CA issuer signs a wildcard certificate for the gateway~\cite{CertManagerCA}. The CA key pair is kept, encrypted, in the repository and created as a Secret by the bootstrap. This keeps the CA the same when an environment is rebuilt (R6), so a device that trusts the CA once keeps trusting the rebuilt environment.

\paragraph{Secrets.} The secrets in Git are encrypted with SOPS, using an age key pair as the encryption key~\cite{SOPS}. SOPS encrypts each value of a YAML file separately and leaves the key names in plain text, so the Git history shows which secret changed in each commit, for example that a single new key was added, while the values stay unreadable. Ansible Vault, the alternative built into Ansible, encrypts the file as one block, so every change replaces the whole file in the diff, and it relies on one password shared by everyone. With SOPS, access is granted per person or machine by adding a public key, and revoking it does not require a new shared secret. SOPS is also not tied to Ansible: the same files can be used from the command line and by other tools, and cloud key services can be added later~\cite{SOPS}. The price of SOPS is that the age private key, kept outside the repository on the control machine, must not be lost.

\paragraph{Git, registry and CI.} Argo CD pulls the repository from a Gitea instance that runs on prgdlv-auto as a standalone host service, together with its container registry and CI runner, all of which I set up for this purpose. I deliberately kept Gitea outside Kubernetes. The cluster is managed through GitOps, so its whole configuration comes from the Git server; if the Git server itself ran in the cluster, the cluster would depend on Git for its configuration and Git on the cluster for its existence, a loop in the dependency graph. A fresh cluster could then not be bootstrapped from Git, and a broken cluster would take down the source needed to repair it. Argo CD addresses Gitea by its VPN IP address rather than by name, because pods cannot be relied on to resolve the VPN's host names.

\paragraph{Access.} The VMs can be reached directly over SSH, which is all that Ansible needs. Checking the platform, however, means opening many services in a browser and from the command line: the demo applications, Argo CD, the Rollouts dashboard, Prometheus, Grafana and Kiali. The NetBird VPN, installed for ease of access, connects the control machine and both VMs in one private network, and its DNS resolves the wildcard names \texttt{*.dev.internal} to the VM. Every service is thus reachable by its own name through the Istio gateway, with a valid certificate, without SSH tunnels or port forwarding for each of them. The same network also keeps the services private (R7).

\paragraph{Delivery controller.} Argo Rollouts provides canary and blue-green releases, traffic mirroring, the experiment step and analysis with several metric providers, all as Kubernetes resources that Argo CD manages like any other~\cite{ArgoRollouts}. Flagger is a comparable progressive delivery operator~\cite{Flagger}; it would replace Argo Rollouts rather than complement it. I chose Argo Rollouts for four reasons, of which the first two were decisive:
\begin{itemize}
	\item \emph{Integration with Argo CD.} Argo Rollouts belongs to the same project as Argo CD, the GitOps controller of the platform~\cite{ArgoCD}, so the whole platform is operated with one family of tools.
	\item \emph{Explicit steps.} A Rollout lists the steps of a release in order: the weights, the pauses, the mirroring, the experiment and the analyses~\cite{ArgoRolloutsCanary}. The evaluation compares strategies that differ in exactly one mechanism, and with Argo Rollouts each strategy is a short Kustomize component that rewrites this list (\autoref{sec:EvalStrategies}). Flagger describes a canary release by parameters, such as the step weight, the maximum weight and the interval of the checks, and derives the steps itself~\cite{FlaggerStrategies}. The order of the mechanisms is then less visible and harder to vary one at a time.
	\item \emph{The experiment step.} An experiment runs fresh pods of the stable and of the new version side by side and compares them~\cite{ArgoRolloutsExperiment}. Flagger has no comparable baseline comparison, so the experiment strategy of the evaluation could not have been built with it.
	\item \emph{Routing kept in Git.} Argo Rollouts only changes the weights of a VirtualService that is defined in the repository~\cite{ArgoRolloutsIstio}. Flagger leaves the Deployment in place but generates further objects itself: a copy of it with the suffix \texttt{-primary}, which serves as the stable version while the original is scaled to zero, and the Services that select the two versions~\cite{FlaggerHowItWorks}. These objects are not in Git, so the state of the cluster would no longer follow from the repository alone.
\end{itemize}
Flagger has two advantages: it needs no new workload resource, and it brings a load tester for workloads that receive no traffic~\cite{FlaggerWebhooks}. The platform covers the second with a Kubernetes Job (\autoref{sec:DesignBlueGreen}), and neither outweighs the reasons above.

\paragraph{Metrics and dashboards.} Why the checks use Prometheus metrics, and not logs or traces, is discussed in \autoref{sec:MetricAnalysis}. Prometheus runs as a standalone server, without Alertmanager, Pushgateway, node exporter and kube-state-metrics, because the checks need none of them and each would cost memory (R5). It scrapes every 10\,s instead of the default minute. With the analysis measuring every 20--30\,s, a one-minute scrape interval would make consecutive measurements read the same data, and a rate over 30\,s would have fewer than two samples; at 10\,s, a one-minute rate window holds about six. Two dashboards complement the Prometheus UI. Grafana shows what happened over a period of time, which is the natural view of a release: the error rate of the canary rising, crossing the threshold, and the traffic returning to the stable version. Kiali draws the live traffic graph of the mesh, which shows the split between the stable and canary Services. Both run without login for viewers, because they are reachable only over the VPN; Kiali is additionally read-only, since any change it made to the routing would be undone by Argo CD within seconds.

%----------------------------------------------------------------------------
\section{Release flow}
\label{sec:DesignReleaseFlow}
%----------------------------------------------------------------------------
\autoref{fig:DesignReleaseFlow} follows one release through the platform. A release starts with a commit to the platform repository that changes the pod template of the application's Rollout. There are two such commits. A code change in the application repository is built by its CI pipeline into an image, pushed to the registry, and the pipeline then commits the new image tag into the platform repository. A configuration change, which is how the faults of the experiments are injected, edits the environment variables of the pod template directly in the platform repository. In both cases Argo CD applies the new Rollout, Argo Rollouts creates a new ReplicaSet and steers traffic to it step by step, and the analysis decides from Prometheus whether to continue.

\begin{figure}[ht]
	\centering
	\resizebox{\textwidth}{!}{%
	\begin{tikzpicture}[
		stage/.style={draw, rounded corners=2pt, align=center, font=\small, text width=2.6cm, minimum height=1.5cm, fill=black!5},
		outcome/.style={draw, rounded corners=2pt, align=center, font=\small, text width=3.6cm, minimum height=1.1cm},
		arr/.style={-{Stealth[length=2mm]}, semithick},
		lbl/.style={font=\scriptsize, align=center}]
		\node[stage] (commit) at (0,0)    {commit to the platform repository\\(image tag or environment)};
		\node[stage] (sync)   at (3.4,0)  {Argo CD syncs the Rollout};
		\node[stage] (rs)     at (6.8,0)  {Argo Rollouts creates a new ReplicaSet};
		\node[stage] (route)  at (10.2,0) {steps set mirror, experiment and route weights};
		\node[stage] (ana)    at (13.6,0) {AnalysisRuns query Prometheus};
		\node[outcome]  (prom)   at (11.9,-2.6) {\textbf{all checks pass:}\\promoted to 100\,\%};
		\node[outcome]  (abort)  at (16.0,-2.6) {\textbf{a check fails:}\\aborted, users back to stable};
		\node[outcome]  (ci)     at (0,2.4)  {code change: CI builds and pushes the image, commits its tag};
		\draw[arr] (ci) -- (commit);
		\draw[arr] (commit) -- (sync);
		\draw[arr] (sync) -- (rs);
		\draw[arr] (rs) -- (route);
		\draw[arr] (route) -- (ana);
		\draw[arr] (ana.south) -- ++(0,-0.6) -| (prom.north);
		\draw[arr] (ana.south) -- ++(0,-0.6) -| (abort.north);
		\draw[arr] (ana.north) -- ++(0,0.6) -| node[lbl, above, pos=0.25]{next step while the checks pass} (route.north);
	\end{tikzpicture}}
	\caption{The flow of a release. A configuration change starts directly with the commit; a code change first goes through the CI pipeline of the application repository.}
	\label{fig:DesignReleaseFlow}
\end{figure}

\paragraph{The image tag is written into the platform repository.} The CI pipeline of the application commits the tag of each new image into the platform repository, instead of the application repository keeping its own manifests with a second Argo CD Application. The platform repository then holds the complete desired state of the cluster in one place, and every release, whatever triggered it, is one commit in one history. Recovering from a faulty release is the revert of that commit, which brings back the previous version at once, without going through the canary steps again, because the previous ReplicaSet is still the stable one.

\paragraph{Faults are configuration.} The faulty versions used in the experiments are not separate builds. The application reads its fault settings from environment variables, and a fault is a set of those variables in the pod template, kept in a small patch file of the platform repository. Argo Rollouts starts a new revision on any change of the pod template, not only on a new image, so a changed variable goes through exactly the same steps and checks as a new image would (\autoref{sec:EvalEnvFaults}). The variables are set in the pod template itself and not taken from a ConfigMap, because a changed ConfigMap does not change the pod template and therefore would not start a release at all. For the same reason every release, including a healthy one, sets its own version label, so that it always changes the template. The first experiments still used one image per fault; the change and its validation are described in \autoref{sec:EvalEnvFaults}.

\paragraph{An aborted release stays visible.} After an abort, Git still describes the faulty version, so the Rollout stays in the \emph{Degraded} state and no new attempt is made until Git changes. This is intended: the platform rolls the users back immediately, but leaves the decision about the faulty commit, a revert or a fix, to a human.

%----------------------------------------------------------------------------
\section{Detection mechanisms}
\label{sec:DesignMechanisms}
%----------------------------------------------------------------------------
The platform has three mechanisms for its canary release and offers blue-green as a separate strategy. The default canary release combines two of them, mirroring and the canary steps; \autoref{fig:DesignTimeline} shows how they follow each other and how many users each exposes to the new version. Mirroring goes first, because it exposes no user at all: a release that fails on copied traffic is stopped before anyone is served by it. It cannot see slow faults, which the canary steps after it catch. The third mechanism, the baseline experiment, is the only relative check; it catches the same faults as the canary steps and costs the most time and pods, so it is not part of the default and is available as a strategy of its own. The default blue-green release likewise has a check before the switch, on generated load, and one after it. All mechanisms can be switched on and off independently, which the experiments use to measure each of them alone and in combination; the first experiments, which tried every combination, led to this choice of defaults (\autoref{sec:EvalStrategies}).

\begin{figure}[ht]
	\centering
	\resizebox{\textwidth}{!}{%
	\begin{tikzpicture}[
		seg/.style={draw, align=center, font=\small, minimum height=1.0cm, inner sep=2pt},
		chk/.style={font=\scriptsize, align=center},
		x=2.4cm]
		% canary row
		\node[font=\small\bfseries, anchor=east] at (-0.2,0) {Canary};
		\node[seg, fill=black!5, minimum width=2.88cm, anchor=west] at (0,0)   {mirror\\0\,\%};
		\node[seg, fill=black!25, minimum width=2.4cm, anchor=west] at (1.2,0) {10\,\%};
		\node[seg, fill=black!35, minimum width=2.4cm, anchor=west] at (2.2,0) {50\,\%};
		\node[seg, fill=black!50, text=white, minimum width=2.4cm, anchor=west] at (3.2,0) {100\,\%};
		\draw[thick] (0,-0.75) -- (1.2,-0.75) node[chk, midway, below]{\texttt{mirror-health}};
		\draw[thick] (1.2,-0.75) -- (3.2,-0.75) node[chk, midway, below]{\texttt{canary-health} (background)};
		% blue-green row
		\node[font=\small\bfseries, anchor=east] at (-0.2,-2.4) {Blue-green};
		\node[seg, fill=black!5, minimum width=2.95cm, anchor=west] at (0,-2.4) {preview\\0\,\%};
		\node[seg, fill=black!50, text=white, minimum width=6.0cm, anchor=west] at (1.23,-2.4) {100\,\%, old version kept running};
		\draw[thick] (0,-3.15) -- (1.23,-3.15) node[chk, midway, below]{load Job, \texttt{preview-load}};
		\draw[thick] (1.23,-3.15) -- (3.73,-3.15) node[chk, midway, below]{\texttt{bg-post-promotion}};
		% time axis
		\draw[-{Stealth[length=2mm]}] (0,-4.2) -- (4.4,-4.2) node[right, font=\scriptsize]{min};
		\foreach \t in {0,1,...,4} \draw (\t,-4.12) -- (\t,-4.28) node[below, font=\scriptsize]{\t};
	\end{tikzpicture}}
	\caption{The default canary release and the default blue-green release in the environment, with the share of users served by the new version in each phase and the check that runs during it. Lengths are approximate.}
	\label{fig:DesignTimeline}
\end{figure}

%----------------------------------------------------------------------------
\subsection{Common design of the checks}
\label{sec:DesignChecks}
%----------------------------------------------------------------------------
All checks are Prometheus queries in AnalysisTemplates, evaluated periodically by Argo Rollouts (\autoref{sec:ArgoRollouts}). They share a set of design decisions, each made in response to a concrete way a check could go wrong.

\paragraph{Whose view of a request.} Each request in the mesh is recorded twice: by the caller's proxy (\texttt{reporter="source"}) and by the receiving pod's proxy (\texttt{reporter="destination"}). The checks of the canary steps and of the experiment use the caller's view, here the ingress gateway, because the gateway labels each request with the Service it routed the request to. The receiving proxy cannot tell \texttt{demo-stable} from \texttt{demo-canary} when both Services select the same pods, which is the normal state between releases. The other checks need the receiving view instead. During mirroring, the gateway does not report the copies it sends, so only the canary's own proxy sees them. Before a blue-green switch, the requests come from the load Job, not through the gateway. After the switch, all requests go to the same active Service, so the requests to the new version can only be separated by the pod that served them. In both cases the new pods are selected by their pod-template hash, which Prometheus copies from the pod labels onto every series of the proxy; Argo Rollouts passes the hash of the new ReplicaSet to the template as an argument.

\paragraph{One value per measurement.} Every query reduces the last minute to a single number: a ratio of request rates for the error share, a 95th percentile from the latency histogram, and the slope of a least-squares line through the container's memory for memory growth. Each measurement therefore summarises about a minute of requests, so a single slow or failed request cannot fail it. The window is one minute because a shorter one would hold too few requests at the traffic of the experiments, and a longer one would react to a fault more slowly.

\paragraph{No data is a failure.} A check must never pass because it saw nothing, for example because the routing is wrong or the canary has not received traffic yet. Every condition therefore starts with \texttt{len(result) > 0}, so that an empty result counts as a failed measurement, and the numerator of each error ratio ends with \texttt{or vector(0)}, so that a version that fails every request yields 0 instead of an empty result. To keep the guard from failing healthy releases, the first measurement waits 30\,s, until Prometheus has at least two samples of the new version in the window.

\paragraph{How many failures abort.} Most checks tolerate one failed measurement and abort on the second (\texttt{failureLimit: 1}), so a single unlucky measurement cannot roll back a healthy release, while a persistent fault is still caught within two measurement intervals. Two kinds of checks abort on the first failure. The memory checks do, because a leak and a healthy pod are far apart (a healthy pod grows by about 0, the leak of the experiments by more than six times the limit), and because the canary lasts only about two minutes in development, too short to wait for a second failure. The success rate during mirroring does too, because no user is exposed at that point: a false alarm costs only a repeated release, while a missed fault reaches real users in the next phase.

\paragraph{Thresholds.} The absolute thresholds are a success rate of at least 95\,\%, a 95th percentile latency of at most 300\,ms and a memory growth of at most 0.05\,MiB/s (3\,MiB per minute). Each lies far from the values of a healthy version (a success rate of 100\,\%, a latency of a few milliseconds and no growth), so that the noise of a healthy release cannot cross it. The experiments later showed that the latency threshold lies inside a bucket of Istio's latency histogram, which makes its estimate coarse; \autoref{sec:EvalMetricValidation} quantifies this and recommends a threshold on a bucket boundary.

%----------------------------------------------------------------------------
\subsection{Canary steps}
\label{sec:DesignCanary}
%----------------------------------------------------------------------------
The canary steps send 10\,\% of the users to the new version, then 50\,\%, and then promote it. Each step is followed by a pause, one minute each, where the longer pauses give the checks more real traffic to judge and more time for slowly developing faults to show. The background analysis \texttt{canary-health} starts at the pause after the first weight step and runs until the release ends. It checks the success rate, the latency and the memory growth of the canary with the absolute thresholds above, every 30\,s.

The memory check needed its own timing. A new pod's memory jumps by about 8\,MiB in its first seconds and then stays flat. A least-squares slope over a window that still contains this jump reads as growth, which aborted a healthy release in a pilot run. The check therefore uses a one-minute window and starts only after 90\,s, when the start-up jump has left the window.

Independently of the steps, every Rollout has a progress deadline of 120\,s with \texttt{progressDeadlineAbort} enabled~\cite{ArgoRolloutsSpec}. A version whose pods never become ready receives no traffic and therefore produces no metrics; without the deadline the release would stay in the \emph{Progressing} state forever, and with it the release is aborted after two minutes. The deadline is a property of the Rollout, not a step, so it protects every strategy.

%----------------------------------------------------------------------------
\subsection{Traffic mirroring}
\label{sec:DesignMirror}
%----------------------------------------------------------------------------
I added mirroring in front of the canary steps and kept those unchanged, so that a fault can be caught on copied traffic before any user is routed to the new version. Istio copies requests to a second destination and discards the responses of the copy, which never reach the client~\cite{IstioMirroring}; Argo Rollouts creates such a mirror route with its \texttt{setMirrorRoute} step, in a route that it manages next to the weighted one~\cite{ArgoRolloutsTraffic}.

The mirror phase consists of five steps. Because the traffic weight of the canary is still 0 during mirroring, and Argo Rollouts sizes the canary from its weight, the canary would have no pod to receive the copies; a \texttt{setCanaryScale} step therefore gives it one pod explicitly, and a second one hands its size back to the weight afterwards~\cite{ArgoRolloutsCanary}. Between them, the mirror route is created, the analysis step \texttt{mirror-health} runs, and the route is removed. Only \texttt{GET} requests to the application's work endpoint are mirrored, because a copied request is executed for real: a non-idempotent request, such as one that writes data, would be performed twice.

\texttt{mirror-health} checks the success rate and the latency of the copies, as recorded by the canary's own proxy (\autoref{sec:DesignChecks}). It takes three measurements 20\,s apart after an initial delay of 30\,s, so the mirror phase lasts about 70\,s. Being an analysis \emph{step} with a fixed number of measurements, it blocks the release until it has finished, and a failure aborts the release before the first user is routed to the new version. It has no memory check: 70\,s is too short for a leak to separate from the start-up jump of memory, and leaks are left to the later mechanisms.

%----------------------------------------------------------------------------
\subsection{Baseline experiment}
\label{sec:DesignExperiment}
%----------------------------------------------------------------------------
The checks above are \emph{absolute}: they compare the new version with fixed thresholds. The experiment step adds a \emph{relative} check, which compares the new version with the old one under the same conditions. It starts two fresh pods side by side for three minutes, a baseline running the stable pod template and a canary running the new one, each behind its own Service and each receiving the same share of the users~\cite{ArgoRolloutsExperiment}. The baseline is started fresh, rather than compared with the long-running stable pods, so that effects of starting up, such as the initial memory allocation and cold caches, and effects of the node, such as CPU pressure, affect both pods equally and cancel out in the difference.

The AnalysisTemplate \texttt{canary-vs-baseline} checks three differences every 20\,s: of the 5xx share (at most 5 percentage points), of the 95th percentile latency (at most 100\,ms) and of the memory growth (at most 0.05\,MiB/s). Latency is compared as a difference, not as a ratio, because the application answers in a few milliseconds, and a ratio would report 7\,ms against 4\,ms as a 75\,\% regression. The memory series come from the kubelet's cAdvisor, which carries no pod labels, and without kube-state-metrics there is no way to join them with the pod-template hash; the memory check therefore selects the two pods by the names the experiment gives them.

Each experiment pod receives 10\,\% of the users, the same share as the first canary step. In real use every request comes from a real user, and a larger experiment weight would show the new version to more users than the first canary step would.

The experiment is not part of the default release. It catches the same faults as the canary steps, so the canary after it would only add time, and it is the most expensive mechanism, at about three minutes and two extra pods per release (\autoref{sec:EvalStrategies}). It remains available as a strategy of its own, the relative alternative to the canary steps.

%----------------------------------------------------------------------------
\subsection{Blue-green release}
\label{sec:DesignBlueGreen}
%----------------------------------------------------------------------------
A Rollout has exactly one strategy, so blue-green is a second copy of the demo application, \texttt{demo-bg}, in its own namespace, running the same image with the same faults~\cite{ArgoRolloutsBlueGreen}. The new version starts behind a preview Service that no user reaches. A pre-promotion analysis then has to pass, after which the active Service is switched to the new version, so all users move at once. A post-promotion analysis watches the new version under the full user traffic and switches the users back if it fails.

The design choices follow from this sequence. The blue-green Rollout uses no traffic routing in Istio: its VirtualServices always point at the active and the preview Service, and a release only changes the selectors of these Services. Argo CD therefore needs no exception for routes that the controller rewrites, unlike for the canary. Promotion is automatic as soon as the pre-promotion analysis passes. The old ReplicaSet is kept for 300\,s after the switch, longer than the post-promotion analysis runs (about 2.5 minutes), so that a rollback is an instant switch of the selector instead of a cold start of the old version. The same progress deadline as for the canary catches a version that never becomes ready. The image of \texttt{demo-bg} is not updated by the CI pipeline: if it were, every build would start a canary and a blue-green release at the same time on the same node, and their metrics and load would mix.

\paragraph{Pre-promotion: metrics on generated load.} Since the preview receives no user traffic, its check needs generated requests. The blue-green strategy has no traffic routing, so it cannot mirror the users' requests to the preview, as the canary does (\autoref{sec:EvalStrategies}). The pre-promotion analysis \texttt{preview-load} therefore starts a Kubernetes Job through the Job provider of Argo Rollouts~\cite{ArgoRolloutsJob}, which sends the same request as the load generator, at the same pace, to the preview Service for 60\,s. The Job only produces load and never fails because of an answer. Two metrics of the same analysis, which run in parallel with the Job, judge the preview pods with the queries, thresholds and timing of \texttt{mirror-health}, so that the check before the switch is the same in both release types. This is the pattern of Flagger's load tester, which generates traffic for the analysis of a workload that receives none~\cite{FlaggerWebhooks}. The Job runs inside the mesh, because strict mutual TLS rejects callers without a sidecar. This requires the native sidecar mode of Kubernetes, in which the proxy runs as an init container that the kubelet stops when the main container exits~\cite{K8sSidecar}; with a classic sidecar, the proxy would keep running and the Job would never complete.

\paragraph{Post-promotion analysis.} \texttt{bg-post-promotion} applies the same three absolute checks as \texttt{canary-health} to the requests served by the new pods, selected by their hash at the receiving proxy, so that the old version's traffic from before the switch does not dilute the one-minute window. Its memory check starts after 90\,s, as in \texttt{canary-health}: without a pre-promotion analysis, the new pod is switched to the users seconds after it starts, while its start-up jump is still in the window. It ends after a fixed number of measurements, about 2.5 minutes after the switch.

%----------------------------------------------------------------------------
\section{Designs considered and not built}
\label{sec:DesignRejected}
%----------------------------------------------------------------------------
\paragraph{Mirroring to the experiment pods.} Mirrored copies could be sent to both experiment pods instead of real users, which would combine the relative check with zero exposure, a design known as a dark launch with a baseline. Diffy works this way: it sends each request to the candidate and to two instances of the old version, and treats differences between the two old instances as noise~\cite{Diffy}. Istio could do this with several mirror destinations, but Argo Rollouts' mirror step can only target the canary Service, and the experiment step has no mirroring option. The route would have to be written into the VirtualService by hand, while Argo CD is configured to leave its routes to Argo Rollouts (\autoref{sec:ImplGitOps}), so such a route would never be applied. Because of this cost, I describe it here but did not build it.

\paragraph{A relative mirror check.} A cheaper relative variant was also considered: during 100\,\% mirroring, the stable version handles exactly the requests that are copied to the canary, so the mirror check could compare the canary's error share with the stable version's instead of with a fixed threshold. This would give a relative error check with no user exposed, but not a relative memory check, because the stable pods are warm and the canary is fresh. It is left for future work.

\paragraph{Statistical and percentage-based comparison.} Statistical canary analysis, as done by Kayenta with a significance test between canary and baseline~\cite{Kayenta}, and a simpler comparison of the canary with the stored metrics of the last stable release were both discussed. I adopted neither: the experiment step already provides a comparison with the old version, and the fixed thresholds were enough for the faults under study.

\paragraph{Other tools for the experiments.} A chaos engineering tool such as Chaos Mesh for environment faults, a load generator such as k6 that records every request on the client side, and a workflow engine such as Argo Workflows to run the trials inside the cluster were considered; I rejected them in favour of the tools already in place. Environment faults were outside the scope (\autoref{sec:DesignRequirements}); the demo application's own fault injection and the curl-based load generator were sufficient; and a workflow engine would have required a new component, a Git token in the cluster and an artifact store, while the actual work of the runner, observing and measuring a release, would have had to be written in any case. The trial runner is therefore a Python program using only the standard library (\autoref{sec:EvalRunner}).
